\newpage
\section*{Part B: Exercise 10.1}

%---------------------------------------------------------------
\subsection*{Exercise 10.1, Question 1}
\textbf{Question.} How many tangents can a circle have?

\textbf{Solution.} At every point of a circle there is one and only one tangent: it is the line drawn through that point perpendicular to the radius to that point (the tangent is perpendicular to the radius through the point of contact). A circle has infinitely many points, so it has infinitely many tangents.

\ans{A circle can have infinitely many tangents.}

%---------------------------------------------------------------
\subsection*{Exercise 10.1, Question 2 (i)}
\textbf{Question.} Fill in the blank: A tangent to a circle intersects it in \underline{\hspace{2cm}} point(s).

\textbf{Solution.} A tangent touches the circle at exactly one point, the point of contact. Every other point of the tangent lies outside the circle.
\slip{A line that meets the circle in \emph{two} points is not a tangent; it is a secant (see part (ii)).}

\ans{one.}

\subsection*{Exercise 10.1, Question 2 (ii)}
\textbf{Question.} Fill in the blank: A line intersecting a circle in two points is called a \underline{\hspace{2cm}}.

\textbf{Solution.} By definition, a line that cuts a circle at two distinct points is a secant of the circle.

\ans{secant.}

\subsection*{Exercise 10.1, Question 2 (iii)}
\textbf{Question.} Fill in the blank: A circle can have \underline{\hspace{2cm}} parallel tangents at the most.

\textbf{Solution.} The tangents at the two ends of a diameter are parallel (both are perpendicular to that diameter). A tangent parallel to these must also be perpendicular to the same diameter at its point of contact, and the only such points on the circle are the two ends of the diameter. So there is no third parallel tangent.
\slip{Having \emph{many} tangents (part of Question 1) is different from having many \emph{parallel} tangents. For one fixed direction there are exactly two tangents, one on each side of the circle.}

\ans{two.}

\subsection*{Exercise 10.1, Question 2 (iv)}
\textbf{Question.} Fill in the blank: The common point of a tangent to a circle and the circle is called \underline{\hspace{2cm}}.

\textbf{Solution.} This is the definition of the point of contact.

\ans{point of contact.}

%---------------------------------------------------------------
\subsection*{Exercise 10.1, Question 3}
\textbf{Question.} A tangent $PQ$ at a point $P$ of a circle of radius $5$ cm meets a line through the centre $O$ at a point $Q$ so that $OQ=12$ cm. Length $PQ$ is:\\
(A) $12$ cm \qquad (B) $13$ cm \qquad (C) $8.5$ cm \qquad (D) $\sqrt{119}$ cm.

\begin{center}
\begin{tikzpicture}[scale=0.38]
\coordinate (O) at (0,0); \coordinate (Q) at (12,0);
\coordinate (P) at ({25/12},{5*sqrt(119)/12});
\draw (O) circle (5);
\fill[blue!8] (O) -- (P) -- (Q) -- cycle;
\draw (P) -- (Q); \draw (O) -- (P); \draw (O) -- (Q);
\draw pic[draw,angle radius=0.45cm] {right angle=O--P--Q};
\dotat{O}{left}\dotat{P}{above}\dotat{Q}{right}
\node[font=\small,below=2pt] at (8.5,0) {$12$ cm};
\node[font=\small,above left] at (1.1,2.3) {$5$ cm};
\end{tikzpicture}%
\figcap{Circle with centre $O$, tangent $PQ$ at $P$, $OQ=12$ cm. The shaded $\tri OPQ$ has a right angle at $P$ (tangent $\perp$ radius).}
\end{center}

\textbf{Solution.} Join $OP$. The tangent is perpendicular to the radius through the point of contact, so $OP\perp PQ$, that is $\ang OPQ=90\deg$. Hence $\tri OPQ$ is a right-angled triangle with the right angle at $P$. Its hypotenuse is the side opposite the right angle, which is $OQ$; the two legs are $OP=5$ cm and $PQ$.

By the Pythagoras theorem, $(\text{hypotenuse})^2=(\text{one leg})^2+(\text{other leg})^2$, so $OQ^2=OP^2+PQ^2$:
\begin{align*}
PQ^2&=OQ^2-OP^2\\
&=12^2-5^2\\
&=144-25=119\\
PQ&=\sqrt{119}\text{ cm}.
\end{align*}
\slip{The right angle is at $P$, not at $O$, so the hypotenuse is $OQ=12$. The squares of the two legs add up to $12^2$, so we \emph{subtract}: $PQ^2=12^2-5^2$. Adding, $12^2+5^2=169$, gives $13$, which is option (B), a tempting wrong choice.}

\textbf{Checking the other options.} $PQ$ is a leg of the right triangle, so $PQ<OQ=12$ cm; this rules out (A) and (B). For (C), $8.5^2=72.25\neq119$.

\textbf{Verification.} Put $PQ^2=119$ back: $OP^2+PQ^2=25+119=144=12^2=OQ^2$.

\ans{Option (D): $PQ=\sqrt{119}$ cm.}

%---------------------------------------------------------------
\subsection*{Exercise 10.1, Question 4}
\textbf{Question.} Draw a circle and two lines parallel to a given line such that one is a tangent and the other, a secant to the circle.

\textbf{Solution.} This is a construction question. Unlike the other questions in this chapter, the figure here \emph{is} the answer, so it shows the finished construction.

\textbf{Steps of construction.}
\begin{enumerate}[label=\arabic*.,leftmargin=*,itemsep=1pt]
\item Draw the given line $\ell$. Draw a circle with centre $O$ and radius $3$ cm.
\item Through $O$ draw the line $m$ perpendicular to $\ell$ (using a set square, or the compass construction for a perpendicular from a point to a line). Let $m$ meet the circle at $P$.
\item At $P$ draw the line $t$ perpendicular to $OP$.
\item Mark a point $R$ on $OP$ with $OR=1.5$ cm. Through $R$ draw the line $s$ perpendicular to $m$.
\end{enumerate}

\begin{center}
\begin{tikzpicture}[scale=0.7]
\begin{scope}[rotate=20]
\coordinate (O) at (0,0); \coordinate (P) at (0,3); \coordinate (R) at (0,1.5);
\draw (O) circle (3);
\draw (-4.5,-4.5) -- (4.5,-4.5);
\draw (-4.5,3) -- (4.5,3);
\draw (-4.5,1.5) -- (4.5,1.5);
\draw[dashed] (0,-5.2) -- (0,3.9);
\coordinate (lL) at (4.5,-4.5); \coordinate (lT) at (4.5,3);
\coordinate (lS) at (4.5,1.5); \coordinate (lM) at (0,3.9);
\end{scope}
\node[right,font=\small] at (lL) {$\ell$ (given line)};
\node[right,font=\small] at (lT) {$t$ (tangent)};
\node[right,font=\small] at (lS) {$s$ (secant)};
\node[above,font=\small] at (lM) {$m$};
\dotat{O}{below right}
\fill (P) circle (1.6pt); \node[font=\small] at ($(P)+(155:0.45)$) {$P$};
\fill (R) circle (1.6pt); \node[font=\small] at ($(R)+(155:0.45)$) {$R$};
\end{tikzpicture}%
\figcap{The tangent $t$ and the secant $s$, both parallel to the given line $\ell$.}
\end{center}

\textbf{Reasoning.}
\begin{itemize}[leftmargin=*,itemsep=1pt]
\item $\ell$, $t$ and $s$ are all perpendicular to the same line $m$, so they are parallel to one another.
\item $t$ is perpendicular to the radius $OP$ at its end $P$ on the circle. Every other point of $t$ is the far end of a hypotenuse from $O$, so it is farther than $OP$ from $O$ and lies outside the circle. So $t$ meets the circle only at $P$: it is a tangent.
\item $OR=1.5$ cm is less than the radius $3$ cm, so $R$ is inside the circle. Any line through an interior point cuts the circle at two points, so $s$ is a secant.
\end{itemize}
\slip{The tangent must be perpendicular to the \emph{radius} $OP$ at $P$, and it must pass through $P$ on the circle. A line through $P$ at any other angle to $OP$ would cut the circle again and be a secant.}

\ans{The line $t$ (tangent) and the line $s$ (secant) are both parallel to the given line $\ell$.}
